% Generated draft: requires mathematical and source review.
\documentclass[11pt]{article}
\usepackage[margin=25mm]{geometry}
\usepackage{fontspec}
\setmainfont{Latin Modern Roman}
\usepackage{amsmath,amssymb,mathtools,unicode-math}
\setmathfont{Latin Modern Math}
\usepackage{enumitem,needspace,xurl,hyperref}
\setlength{\parindent}{0pt}
\setlength{\parskip}{5pt}
\setlength{\emergencystretch}{3em}
\newcommand{\N}{\mathbb{N}}
\newcommand{\Z}{\mathbb{Z}}
\newcommand{\Q}{\mathbb{Q}}
\newcommand{\R}{\mathbb{R}}
\newcommand{\C}{\mathbb{C}}
\newcommand{\F}{\mathbb{F}}
\newcommand{\A}{\mathbb{A}}
\newcommand{\PP}{\mathbb{P}}
\newcommand{\E}{\mathbb{E}}
\newcommand{\eps}{\varepsilon}
\newcommand{\dd}{\,\mathrm{d}}
\begin{document}
\section*{UnsolvedMath ID: 30002587}
\textbf{Source problem code:} catalogue entry
\subsection{Definitions and mathematical statement}
Let \(n\geq 1\) be an integer and let \(\mathbb B_n=\{z\in\C^n:\|z\|_2<1\}\). The Drury--Arveson space \(H_n^2\) is the reproducing-kernel Hilbert space of holomorphic functions on \(\mathbb B_n\) with kernel
\[
K_z(w)=\frac{1}{1-\langle w,z\rangle},\qquad \langle w,z\rangle=\sum_{j=1}^n w_j\overline{z_j}.
\]
A holomorphic function \(f\) is a multiplier of \(H_n^2\) if multiplication defines a bounded operator \(M_f:H_n^2\to H_n^2\). Equivalently, there exists \(C<\infty\) such that, for every \(h\in H_n^2\), one has \(fh\in H_n^2\) and \(\|fh\|_{H_n^2}\leq C\|h\|_{H_n^2}\).

The source retains the question whether, for every \(f\in H_n^2\),
\[
\sup_{z\in\mathbb B_n}\|f k_z\|_{H_n^2}<\infty
\quad\Longrightarrow\quad
f\in\operatorname{Mult}(H_n^2).
\]
Here finiteness includes the requirement that \(f k_z\in H_n^2\) for every \(z\in\mathbb B_n\). The supplied record does not define \(k_z\), specify its normalization, or specify the intended dimension or range of dimensions. In particular, the unnormalized kernel \(K_z\) and the normalized kernel
\[
\widehat K_z(w)=\frac{K_z(w)}{\sqrt{K_z(z)}}=\frac{\sqrt{1-\|z\|_2^2}}{1-\langle w,z\rangle}
\]
are distinct conventions. If \(k_z=\widehat K_z\) is intended, the displayed hypothesis becomes
\[
\forall z\in\mathbb B_n:\ f\widehat K_z\in H_n^2,
\qquad
\sup_{z\in\mathbb B_n}\|f\widehat K_z\|_{H_n^2}<\infty.
\]
This is a conditional interpretation, not a verified recovery of the source notation.

\emph{Source qualification.} The supplied statement-verification field says that the official report immediately answers the retained question negatively by exhibiting a function satisfying the displayed bound that is not a multiplier. The counterexample and its dimensional and normalization hypotheses are not supplied. The record therefore labels the statement \texttt{unrecoverable} and promotes no separate clean statement. Its supplied status is \texttt{partially\_solved}. Separately, the supplied dated literature assessment concerns a simpler direct criterion for special derivative measures: it says that a full Carleson-measure characterization gives a solution in principle, that the requested simpler criterion has not been obtained, and that natural kernel tests are insufficient. No derivative measures or precise criterion are defined in this record, so that discussion does not determine an additional mathematical statement.
\subsection{Short English statement}
The retained question asks whether uniform bounds on products of a Drury--Arveson function with kernel functions force bounded multiplication on the whole space. The supplied verification says the report answers this negatively, but the kernel normalization, relevant dimensions, and counterexample are missing. The separate derivative-measure discussion has no supplied precise formulation. This is a statement-only draft requiring human mathematical review.
\subsection{Sources}
\begin{itemize}
\item Supplied local UnsolvedMath registry record, permanent ID \(30002587\), source problem code \texttt{catalogue entry}, title ``Answered Drury--Arveson Kernel-Norm Question,'' in the Oberwolfach Reports Open Problems collection. No individual proposer is identified; the supplied proposed year is 2014.
\item Supplied report citation: \emph{Hilbert Modules and Complex Geometry} (2015), Oberwolfach Reports, pp.~1139--1219, DOI \texttt{10.4171/owr/2014/21}. \url{https://doi.org/10.4171/owr/2014/21}. This is also the supplied literature source labeled ``Oberwolfach report.'' The supplied citation year is retained alongside the 2014 DOI component and proposed year without resolving their relationship.
\item Supplied literature source labeled ``Carleson-measure characterization'': \url{https://arxiv.org/abs/1310.4820}. The record supplies no authors, title, or theorem locator.
\item Supplied literature source labeled ``Multiplier counterexample'': \url{https://doi.org/10.1512/IUMJ.2015.64.5506}. The record supplies no authors, title, or theorem locator.
\item Literature assessment checked at \texttt{2026-08-21}: supplied dated curation, not independently verified. It labels the problem partially solved and states that the full Carleson-measure characterization supplies a solution in principle, while a simpler direct criterion for the special derivative measures has not been obtained and natural kernel tests are insufficient. The background explicitly calls this dated literature triage derived from supplied OWR curation and requests independent verification of time-sensitive status and cited claims.
\end{itemize}
\end{document}
